\documentclass[10pt,a4paper]{amsart}
\usepackage{amsmath,amssymb,amsthm,graphicx}
% Retained preamble for the Stacks Project chapter {CH} (shared project preamble).
% The chapter file pinned-original-source.tex only carries \input{{preamble}} on its first line,
% so no preamble text lives inside the chapter itself.  The Stacks Project uses one single shared
% preamble for every chapter: the snapshot stages it three times (sources/stacks-chow/preamble.tex,
% sources/stacks-etale-cohomology/preamble.tex and focused-seeds/AGP/sources/stacks-snapshot-20261011/
% preamble.tex) and those copies are identical in size.  The block below is copied mechanically,
% line for line, from sources/stacks-chow/preamble.tex lines 14-18 and 154-217: the Xy-pic setup and
% the native theorem declarations together with the mathematical macros.  Nothing was added, removed
% or redefined; the wrapper class and the amsmath/amssymb/amsthm/graphicx base packages are supplied
% by the assembler.
% For commutative diagrams we use Xy-pic
% Chapter: stacks-spaces-more-morphisms.tex -- More on Morphisms of Spaces.
\usepackage[all]{xy}

% We use 2cell for 2-commutative diagrams.
\xyoption{2cell}

% Theorem environments.
%
\theoremstyle{plain}
\newtheorem{theorem}[subsection]{Theorem}
\newtheorem{proposition}[subsection]{Proposition}
\newtheorem{lemma}[subsection]{Lemma}

\theoremstyle{definition}
\newtheorem{definition}[subsection]{Definition}
\newtheorem{example}[subsection]{Example}
\newtheorem{exercise}[subsection]{Exercise}
\newtheorem{situation}[subsection]{Situation}

\theoremstyle{remark}
\newtheorem{remark}[subsection]{Remark}
\newtheorem{remarks}[subsection]{Remarks}

\numberwithin{equation}{subsection}

% Macros
%
\def\lim{\mathop{\mathrm{lim}}\nolimits}
\def\colim{\mathop{\mathrm{colim}}\nolimits}
\def\Spec{\mathop{\mathrm{Spec}}}
\def\Hom{\mathop{\mathrm{Hom}}\nolimits}
\def\Ext{\mathop{\mathrm{Ext}}\nolimits}
\def\SheafHom{\mathop{\mathcal{H}\!\mathit{om}}\nolimits}
\def\SheafExt{\mathop{\mathcal{E}\!\mathit{xt}}\nolimits}
\def\Sch{\mathit{Sch}}
\def\Mor{\mathop{\mathrm{Mor}}\nolimits}
\def\Ob{\mathop{\mathrm{Ob}}\nolimits}
\def\Sh{\mathop{\mathit{Sh}}\nolimits}
\def\NL{\mathop{N\!L}\nolimits}
\def\CH{\mathop{\mathrm{CH}}\nolimits}
\def\proetale{{pro\text{-}\acute{e}tale}}
\def\etale{{\acute{e}tale}}
\def\QCoh{\mathit{QCoh}}
\def\Ker{\mathop{\mathrm{Ker}}}
\def\Im{\mathop{\mathrm{Im}}}
\def\Coker{\mathop{\mathrm{Coker}}}
\def\Coim{\mathop{\mathrm{Coim}}}

% Boxtimes
%
\DeclareMathSymbol{\boxtimes}{\mathbin}{AMSa}{"02}

%
% Macros for moduli stacks/spaces
%
\def\QCohstack{\mathcal{QC}\!\mathit{oh}}
\def\Cohstack{\mathcal{C}\!\mathit{oh}}
\def\Spacesstack{\mathcal{S}\!\mathit{paces}}
\def\Quotfunctor{\mathrm{Quot}}
\def\Hilbfunctor{\mathrm{Hilb}}
\def\Curvesstack{\mathcal{C}\!\mathit{urves}}
\def\Polarizedstack{\mathcal{P}\!\mathit{olarized}}
\def\Complexesstack{\mathcal{C}\!\mathit{omplexes}}
% \Pic is the operator that assigns to X its picard group, usage \Pic(X)
% \Picardstack_{X/B} denotes the Picard stack of X over B
% \Picardfunctor_{X/B} denotes the Picard functor of X over B
\def\Pic{\mathop{\mathrm{Pic}}\nolimits}
\def\Picardstack{\mathcal{P}\!\mathit{ic}}
\def\Picardfunctor{\mathrm{Pic}}
\def\Deformationcategory{\mathcal{D}\!\mathit{ef}}
\newenvironment{reference}{}{}
\newenvironment{slogan}{}{}
\newenvironment{history}{}{}
\usepackage{cleveref}
\crefname{originallemma}{original lemma}{original lemmas}
\Crefname{originallemma}{Original lemma}{Original lemmas}
\makeatletter
\newlabel{lemma-bounded-on-fibres}{{0GFK}{source}{}{sourcepointer-1}{}}
\newlabel{lemma-bounded-on-fibres@cref}{{[originallemma][0][]0GFK}{[1][0][]source}{}{}{}}
\makeatother
\crefname{originallemma}{original lemma}{original lemmas}
\Crefname{originallemma}{Original lemma}{Original lemmas}
\makeatletter
\newlabel{lemma-characterize-pseudo-coh-improved}{{0GES}{source}{}{sourcepointer-2}{}}
\newlabel{lemma-characterize-pseudo-coh-improved@cref}{{[originallemma][0][]0GES}{[1][0][]source}{}{}{}}
\makeatother
\crefname{originallemma}{original lemma}{original lemmas}
\Crefname{originallemma}{Original lemma}{Original lemmas}
\makeatletter
\newlabel{lemma-relative-pseudo-coherent-is-moot}{{0DII}{source}{}{sourcepointer-3}{}}
\newlabel{lemma-relative-pseudo-coherent-is-moot@cref}{{[originallemma][0][]0DII}{[1][0][]source}{}{}{}}
\makeatother
\crefname{originallemma}{original lemma}{original lemmas}
\Crefname{originallemma}{Original lemma}{Original lemmas}
\makeatletter
\newlabel{sites-cohomology-lemma-dual-perfect-complex}{{08JJ}{source}{}{sourcepointer-4}{}}
\newlabel{sites-cohomology-lemma-dual-perfect-complex@cref}{{[originallemma][0][]08JJ}{[1][0][]source}{}{}{}}
\makeatother
\crefname{originallemma}{original lemma}{original lemmas}
\Crefname{originallemma}{Original lemma}{Original lemmas}
\makeatletter
\newlabel{spaces-perfect-lemma-compare-base-change}{{08IR}{source}{}{sourcepointer-5}{}}
\newlabel{spaces-perfect-lemma-compare-base-change@cref}{{[originallemma][0][]08IR}{[1][0][]source}{}{}{}}
\makeatother
\crefname{originallemma}{original lemma}{original lemmas}
\Crefname{originallemma}{Original lemma}{Original lemmas}
\makeatletter
\newlabel{spaces-perfect-proposition-detecting-bounded-below}{{0GFI}{source}{}{sourcepointer-6}{}}
\newlabel{spaces-perfect-proposition-detecting-bounded-below@cref}{{[originallemma][0][]0GFI}{[1][0][]source}{}{}{}}
\makeatother
\crefname{originallemma}{original lemma}{original lemmas}
\Crefname{originallemma}{Original lemma}{Original lemmas}
\makeatletter
\newlabel{spaces-perfect-remark-pullback-generator}{{0E4R}{source}{}{sourcepointer-7}{}}
\newlabel{spaces-perfect-remark-pullback-generator@cref}{{[originallemma][0][]0E4R}{[1][0][]source}{}{}{}}
\makeatother
\crefname{originallemma}{original lemma}{original lemmas}
\Crefname{originallemma}{Original lemma}{Original lemmas}
\makeatletter
\newlabel{spaces-perfect-theorem-bondal-van-den-Bergh}{{09IY}{source}{}{sourcepointer-8}{}}
\newlabel{spaces-perfect-theorem-bondal-van-den-Bergh@cref}{{[originallemma][0][]09IY}{[1][0][]source}{}{}{}}
\makeatother
\title{Relative perfection over a ring tested against all perfect pairings, via fibre boundedness and a generator}
\author{The Stacks Project Authors}
\date{stacks-project:spaces-more-morphisms.tex:local-snapshot-20261011; GNU FDL 1.2 or later}
\begin{document}
\maketitle
\noindent\textit{Editorial note.} Statement span is lines 12780-12787 of sources/stacks-spaces-more-morphisms/pinned-original-source.tex; proof span is lines 12789-12829, the complete author proof (39 body lines).\par\medskip
\noindent\textit{Editorial note.} Navigation stubs: 8 cited-but-not-included labels are supplied as non-load-bearing pointers with their pinned tags where the snapshot tag index has them.\par\medskip
\noindent\textit{Editorial note.} extra\_packages is empty on purpose: the retained preamble already loads xy with options and a second declaration clashes in the assembler wrapper.\par\medskip
\noindent\textit{Editorial note.} Text is retained verbatim, including the author's own typos and placeholders.\par\medskip
\noindent\textit{Editorial note.} Counted B side (both complete). Second load-bearing use of wave-17 cards L16169/L16171.\par\medskip
\noindent\textit{Allowed standard prerequisites (editorial).}\begin{itemize}
\item Explicit prerequisite boundary. Chapter lemmas characterize-pseudo-coh-improved, relative-pseudo-coherent-is-moot (L16171), bounded-on-fibres (L16169); Derived Categories of Spaces: detecting-bounded-below, compare-base-change, bondal-van-den-Bergh, pullback-generator; Cohomology on Sites, lemma dual-perfect-complex.
\end{itemize}\medskip
\ifdefined\hypertarget\hypertarget{sourcepointer-1}{}\fi
\noindent\textit{Source reference.} Original lemma 0GFK: sources/stacks-spaces-more-morphisms/pinned-original-source.tex, source lines 12759--12769. Cited by the retained statement or proof but not part of the collected spans; supplied only as a navigation pointer. The argument does not depend on its content.\par\medskip
\ifdefined\hypertarget\hypertarget{sourcepointer-2}{}\fi
\noindent\textit{Source reference.} Original lemma 0GES: sources/stacks-spaces-more-morphisms/pinned-original-source.tex, source lines 12230--12237. Cited by the retained statement or proof but not part of the collected spans; supplied only as a navigation pointer. The argument does not depend on its content.\par\medskip
\ifdefined\hypertarget\hypertarget{sourcepointer-3}{}\fi
\noindent\textit{Source reference.} Original lemma 0DII: sources/stacks-spaces-more-morphisms/pinned-original-source.tex, source lines 10860--10868. Cited by the retained statement or proof but not part of the collected spans; supplied only as a navigation pointer. The argument does not depend on its content.\par\medskip
\ifdefined\hypertarget\hypertarget{sourcepointer-4}{}\fi
\noindent\textit{Source reference.} Original lemma 08JJ: sources/stacks-spaces-more-morphisms/pinned-original-source.tex, source lines 12829--12829. Cited by the retained statement or proof but not part of the collected spans; supplied only as a navigation pointer. The argument does not depend on its content.\par\medskip
\ifdefined\hypertarget\hypertarget{sourcepointer-5}{}\fi
\noindent\textit{Source reference.} Original lemma 08IR: sources/stacks-spaces-more-morphisms/pinned-original-source.tex, source lines 12829--12829. Cited by the retained statement or proof but not part of the collected spans; supplied only as a navigation pointer. The argument does not depend on its content.\par\medskip
\ifdefined\hypertarget\hypertarget{sourcepointer-6}{}\fi
\noindent\textit{Source reference.} Original lemma 0GFI: sources/stacks-spaces-more-morphisms/pinned-original-source.tex, source lines 12829--12829. Cited by the retained statement or proof but not part of the collected spans; supplied only as a navigation pointer. The argument does not depend on its content.\par\medskip
\ifdefined\hypertarget\hypertarget{sourcepointer-7}{}\fi
\noindent\textit{Source reference.} Original lemma 0E4R: sources/stacks-spaces-more-morphisms/pinned-original-source.tex, source lines 12829--12829. Cited by the retained statement or proof but not part of the collected spans; supplied only as a navigation pointer. The argument does not depend on its content.\par\medskip
\ifdefined\hypertarget\hypertarget{sourcepointer-8}{}\fi
\noindent\textit{Source reference.} Original lemma 09IY: sources/stacks-spaces-more-morphisms/pinned-original-source.tex, source lines 12829--12829. Cited by the retained statement or proof but not part of the collected spans; supplied only as a navigation pointer. The argument does not depend on its content.\par\medskip
\begin{lemma}
\label{lemma-characterize-relatively-perfect}
Let $A$ be a ring. Let $X$ be an algebraic space separated, of
finite presentation, and flat over $A$. Let $K \in D_\QCoh(\mathcal{O}_X)$.
If  $R \Gamma (X, E \otimes^\mathbf{L} K)$ is perfect in
$D(A)$ for every perfect $E \in D(\mathcal{O}_X)$, then $K$ is
$\Spec(A)$-perfect.
\end{lemma}

\begin{proof}
By Lemma \ref{lemma-characterize-pseudo-coh-improved},
$K$ is pseudo-coherent relative to $A$. By Lemma
\ref{lemma-relative-pseudo-coherent-is-moot},
$K$ is pseudo-coherent in $D(\mathcal{O}_X)$.
By Derived Categories of Spaces, Proposition
\ref{spaces-perfect-proposition-detecting-bounded-below}
we see that $K$ is in $D^-(\mathcal{O}_X)$.
Let $\mathfrak{p}$ be a prime ideal of $A$ and denote
$i : Y \to X$ the inclusion of the scheme theoretic
fibre over $\mathfrak{p}$, i.e., $Y$ is a scheme over $\kappa(\mathfrak p)$.
By Lemma \ref{lemma-bounded-on-fibres},
we will be done if we can show  $Li^*(K)$ is bounded below. 
Let $G \in D_{perf} (\mathcal{O}_X)$ be a perfect
complex which generates $D_\QCoh (\mathcal{O}_X)$,
see Derived Categories of Spaces, Theorem
\ref{spaces-perfect-theorem-bondal-van-den-Bergh}.
We have
\begin{align*}
R\Hom _{\mathcal{O}_Y}(Li^*(G), Li^*(K))
& =
R\Gamma(Y, Li^*(G ^\vee \otimes ^\mathbf{L} K)) \\
& =
R\Gamma(X, G^\vee \otimes ^{\mathbf{L}} K)
\otimes^\mathbf{L}_A \kappa(\mathfrak{p})
\end{align*}
The first equality uses that $Li^*$ preserves perfect objects and duals
and Cohomology on Sites, Lemma
\ref{sites-cohomology-lemma-dual-perfect-complex}; we omit
some details. The second equality follows from
Derived Categories of Spaces, Lemma
\ref{spaces-perfect-lemma-compare-base-change}
as $X$ is flat over $A$. It follows from our hypothesis that this is a
perfect object of $D(\kappa(\mathfrak{p}))$. The object
$Li^*(G) \in D_{perf}(\mathcal{O}_Y)$  generates $D_\QCoh(\mathcal{O}_Y)$ by
Derived Categories of Spaces, Remark
\ref{spaces-perfect-remark-pullback-generator}.
Hence Derived Categories of Spaces, Proposition
\ref{spaces-perfect-proposition-detecting-bounded-below}
now implies that $Li^*(K)$ is bounded below and we win.
\end{proof}

\end{document}
